% 第33批量子物理高难题候选：双delta势与拓扑Dirac场
\subsection*{Q226\quad 对称双吸引delta势的偶奇束缚态与隧穿劈裂}
一维粒子质量$m$，$V(x)=-g[\delta(x-a)+\delta(x+a)]$，$g>0,a>0$，$\hbar$保留。束缚态能量$E=-\hbar^2\kappa^2/(2m)$，$\kappa>0$。
(1) 由波函数连续和导数跳跃推导偶态、奇态的束缚能方程。
(2) 证明奇态出现需要满足最小势阱间距条件，偶态始终存在。
(3) 在$a\kappa_0\gg1$时求两态能量的指数劈裂，$\kappa_0=mg/\hbar^2$。
(4) 讨论$a\to0$时两delta融合的单势阱极限及态数变化。
\par\textbf{逐问解析}
(1) 无delta处满足$\psi''=\kappa^2\psi$。偶态在$|x|<a$取$A\cosh(\kappa x)$，$x>a$取$Be^{-\kappa(x-a)}$；奇态内区取$A\sinh(\kappa x)$，外区类似衰减。每个delta处$\psi'(a^+)-\psi'(a^-)=-2mg\psi(a)/\hbar^2=-2\kappa_0\psi(a)$。结合连续性，偶态满足$\kappa[1+\tanh(\kappa a)]=2\kappa_0$，等价于$\kappa=\kappa_0[1+e^{-2\kappa a}]$；奇态满足$\kappa[1+\coth(\kappa a)]=2\kappa_0$，等价于$\kappa=\kappa_0[1-e^{-2\kappa a}]$。
(2) 对偶态，右边$\kappa_0(1+e^{-2\kappa a})$随$\kappa$从$2\kappa_0$连续降至$\kappa_0$，左边严格上升，因此有唯一正交点。奇态写成$1=\kappa_0(1-e^{-2a\kappa})/\kappa$，右边从$2a\kappa_0$随$\kappa$严格降为0，只有$2a\kappa_0>1$时有非零$\kappa$解；等号仅有$\kappa\to0$阈值态而不可归一。
(3) 若$a\kappa_0\gg1$，令小参数$\varepsilon=e^{-2a\kappa_0}$，用不动点展开得$\kappa_e=\kappa_0(1+\varepsilon+O(a\kappa_0\varepsilon^2))$、$\kappa_o=\kappa_0(1-\varepsilon+O(a\kappa_0\varepsilon^2))$。所以$E_o-E_e=\hbar^2(\kappa_e^2-\kappa_o^2)/(2m)=2\hbar^2\kappa_0^2\varepsilon/m+O(a\kappa_0\varepsilon^2)$。劈裂指数按两阱间隧穿距离$2a$变化。
(4) 在$a\to0$时分布势趋$-2g\delta(x)$，单delta束缚态衰减常数为$2mg/\hbar^2=2\kappa_0$，恰由偶态方程得到$\kappa_e\to2\kappa_0$；奇态存在条件$2a\kappa_0>1$终将失效，其能量上升至连续谱阈值而消失。此极限同时核对了势强度与态数，不能简单把两个独立单阱能量相加。
\par\medskip
\subsection*{Q227\quad 二维质量Dirac带的Berry曲率与整数量子化正则化}
考虑二维连续两带哈密顿量$H(\boldsymbol k)=v(k_x\sigma_x+k_y\sigma_y)+M(k)\sigma_z$，$v>0$；先取常质量$M(k)=m\ne0$，再用$M(k)=m-B|\boldsymbol k|^2$作高动量正则化，$B\ne0$。以Berry联络$\mathcal A=i\langle u_-|\nabla_{\boldsymbol k}u_-\rangle$定义下能带曲率。
(1) 求两能带及常质量下能带的Berry曲率分布。
(2) 在连续动量平面积分曲率，解释为什么仅得半整数通量。
(3) 对Wilson型$M=m-Bk^2$计算紧化平面的Chern数并指出相变。
(4) 讨论Berry曲率与Hall电导符号、规范选择及晶格正则化之间的关系。
\par\textbf{逐问解析}
(1) 记$\boldsymbol d=(vk_x,vk_y,m)$，$E_\pm=\pm|\boldsymbol d|=\pm\sqrt{v^2k^2+m^2}$。下带本征矢的Bloch球自旋反平行于$\boldsymbol d$，按照题设$\mathcal A=i\langle u|\nabla u\rangle$约定，下带曲率为$\Omega_-(\boldsymbol k)=\tfrac12\hat{\boldsymbol d}\cdot(\partial_{k_x}\hat{\boldsymbol d}\times\partial_{k_y}\hat{\boldsymbol d})=mv^2/[2(m^2+v^2k^2)^{3/2}]$。它在$k=0$附近最强，随$k^{-3}$衰减，质量变号时整体反号。
(2) 定义$\mathcal C=(2\pi)^{-1}\int_{\mathbb R^2}\Omega_-\,d^2k$。极坐标给$\mathcal C=\int_0^\infty [mv^2k/(2(m^2+v^2k^2)^{3/2})]\,dk=\operatorname{sgn}(m)/2$。这不是允许晶格绝缘体的整数Chern数：常质量连续模型在$k\to\infty$方向依赖角度，无法将整个动量无穷远视为同一个紧化点，故半整数反映未完备的紫外正则化。
(3) 对$M(k)=m-Bk^2$，$\hat{\boldsymbol d}$在$k=0$指向$\operatorname{sgn}(m)\hat z$，无穷远指向$-\operatorname{sgn}(B)\hat z$，从而动量平面可紧化为球。对径向纹理，设$\cos\theta(k)=M(k)/\sqrt{v^2k^2+M(k)^2}$，方位角为$\varphi$，下带曲率积分化为$\mathcal C=\frac12[\cos\theta(0)-\cos\theta(\infty)]=[\operatorname{sgn}(m)+\operatorname{sgn}(B)]/2$，结果必为$0,\pm1$。只有$m=0$或需要跨越模型正则化参数$B=0$的间隙/边界改变时才可变换拓扑数。
(4) 对完全填满的孤立带，Hall响应与积分Berry曲率成正比，但$\sigma_{xy}$的总符号取决于电荷符号和所选$\mathcal A$、响应定义的约定，应统一固定后比较。局域规变$|u_-\rangle\to e^{i\chi(\boldsymbol k)}|u_-\rangle$改变$\mathcal A$为$\mathcal A-\nabla\chi$，曲率及$\mathcal C$不变。真实晶格的Brillouin区本身紧致，额外高能带负责完整的整数化，不能把单个无正则化Dirac锥的半整数直接称为独立可测的整数Hall平台。
\par\medskip
